% =============================================================================
% Report - Plantilla limpia basada en REVTeX 4.2
% Sistemas Complejos - Universidad de Granada
% =============================================================================
%
% Compilacion:
%   latexmk -pdf Report.tex      (o el boton Build de LaTeX Workshop)
%
% Para alternar entre dos formatos:
%   reprint  -> doble columna, estilo articulo final (default).
%   preprint -> una columna, mas margen, util para borradores.

\documentclass[
  reprint,
  amsmath,amssymb,
  aps,
  pra,            % journal especifico: elimina el warning "No journal specified"
  %nofootinbib,
  %floatfix,
]{revtex4-2}

% --- Paquetes ----------------------------------------------------------------
\usepackage[utf8]{inputenc}                % acentos en el .tex
\usepackage[spanish,es-tabla]{babel}       % titulos y captions en espanol
\usepackage{graphicx}                      % \includegraphics
\usepackage{bm}                            % negrita en math
\usepackage{dcolumn}                       % alinear columnas por decimal
\usepackage{hyperref}                      % links clicables
\usepackage[capitalise]{cleveref}          % \cref inteligente (carga DESPUES de hyperref)

% Carpeta donde estan las figuras (opcional; ajusta o quita si las pones aqui mismo).
\graphicspath{{figs/}{../resultados/}}

% =============================================================================
% Frontmatter
% =============================================================================
\begin{document}

\title{Modelo de Kuramoto: sincronizaci\'on colectiva en sistemas complejos}

\author{Martín Iturralde Jara}
\author{Rafael García Millán}
 \email{nose@correo.ugr.es}
 \email{rgm220703@correo.ugr.es}
\affiliation{Universidad de Granada, Facultad de Ciencias}

\date{\today}

\begin{abstract}
Resumen del trabajo: una o dos frases con el problema, otra con el m\'etodo
y otra con el resultado principal. Evita ecuaciones y referencias aqu\'i.
\end{abstract}

\maketitle


% =============================================================================
% 1. Introduccion
% =============================================================================
\section{Introducci\'on}
\label{sec:intro}

Contexto del modelo de Kuramoto, su relevancia y motivaci\'on del trabajo.
% Cuando tengas entradas en ReportNotes.bib, cita asi:
%   Cita a~\cite{acebron2005} para introducir el modelo de campo medio.


% =============================================================================
% 2. Modelo y metodos numericos
% =============================================================================
\section{Modelo y m\'etodos}
\label{sec:metodos}

La dinámica viene dada por
\begin{equation}
  \dot{\theta}_i(t) = \omega_i + K \sum_{j=1}^{N} A_{ij}\,\sin(\theta_j - \theta_i),
  \label{eq:kuramoto}
\end{equation}
donde el parámetro de orden complejo se define como
\begin{equation}
  R(t)\,e^{i\Psi(t)} = \frac{1}{N}\sum_{j=1}^{N} e^{i\theta_j(t)}.
  \label{eq:R}
\end{equation}

Detalles de la integración (Euler explícito, parada por convergencia, etc.).


% =============================================================================
% 3. Resultados
% =============================================================================
\section{Resultados}
\label{sec:resultados}

\subsection{Tarea 1: transici\'on cl\'asica de Kuramoto}
\label{sec:tarea1}

% --- Ejemplo de figura. Descomenta y ajusta el path cuando tengas el .png:
% \begin{figure}[ht]
%   \centering
%   \includegraphics[width=\columnwidth]{R_vs_K.png}
%   \caption{$\langle R\rangle$ frente a $K$ para distintos $\sigma$.}
%   \label{fig:tarea1_R}
% \end{figure}
% Referenciala con: en \cref{fig:tarea1_R} se observa\dots

\subsection{Tarea 2: red modular}
\label{sec:tarea2}

\subsection{Tarea 3: estructura jer\'arquica}
\label{sec:tarea3}

\subsection{Tarea 4: hacia redes tipo cerebro}
\label{sec:tarea4}


% =============================================================================
% 4. Discusion y conclusiones
% =============================================================================
\section{Discusi\'on y conclusiones}
\label{sec:conclusiones}


% =============================================================================
% Agradecimientos (opcional)
% =============================================================================
\begin{acknowledgments}
% Texto de agradecimientos, financiacion, etc.
\end{acknowledgments}


% =============================================================================
% Bibliografia
% =============================================================================
% Las entradas viven en ReportNotes.bib. Se compila con bibtex (latexmk lo
% lanza solo). Cita con \cite{clave} o \cref{eq:...} para ecuaciones/figuras.
\bibliography{ReportNotes}

\end{document}
